\section{Conclusions}

In this work we have established the empirical efficacy of the
notion of rich subgroup fairness and the algorithm of
\cite{KNRW18} on four fairness-sensitive datasets, and the necessity of
explicitly enforcing subgroup (as opposed to only marginal) fairness.
There are
a number of interesting directions for further experimental
work we plan to pursue, including:
\begin{itemize}
\item Experiments with richer Learner model classes $\cH$, while
	keeping the Auditor subgroup class $\cF$ relatively simple
	and fixed. One conjecture is that by making the hypothesis
	space richer, more appealing Pareto curves may be achieved.
	There is also some rationale for keeping $\cF$ simple, since
	we would like to have some intuitive interpretation of what
	the subgroups represent, while the same constraint may not
	hold for $\cH$.
\item Implementation and experimentation with the no-regret algorithm
	of \cite{KNRW18}, which may have superior convergence and other
	properties due to its stronger theoretical guarantees.
\item Experiments on the generalization performance of subgroup
	fairness in the form of test-set Pareto curves. While as mentioned,
	standard VC theory can be applied to obtain worst-case bounds, one
	might expect even better empirical generalization.
\end{itemize}



